% Mechanical relocation generated from the preservation ledger.
% Existing prose is included byte-for-byte from preserved blocks.

\section{Advisory Model Relaxations}
\label{sec:advisory-model-relaxations}

Residual and sensitivity diagnostics may suggest a broad-period cost effect, a
destination-zone attraction effect, an origin-zone production effect, or a
transfer-place correction. Such a suggestion is a candidate scientific
hypothesis, not an automatic model-selection decision. A residual pattern can
also be caused by erroneous counts, an incorrect stop mapping, timetable
misalignment, or an inadequate observation distribution.

Every proposed relaxation must state its added parameters, normalization,
required metadata, prior, expected identifying variation, and minimum
observation support. It begins at a zero-effect value that reproduces the parent
model. Only after this nesting check should the child be fitted and compared
using adequacy, grouped holdout, curvature, and sensitivity to initialization
and prior scale.

Retain the smallest model that passes the declared checks. A relaxation that
improves calibration fit but produces unstable parameters, worsens held-out
performance, or merely absorbs a data defect should not be adopted. Automated
scores can rank candidates for review; they cannot replace the scientific
decision.

Relaxations can now be expressed as complete serialized specifications rather
than Python edits. In addition to the three original atomic helpers, the
declarative interface supports origin, destination, time-period, origin--time,
destination--time, zone-pair, and custom-group effects, plus an optional smooth
time basis. The validation summary lists the generated parameter blocks,
required feature mappings, calibration and excluded rows, regularization,
feature-cache and structural-zero fingerprints, and all identifiability
warnings before optimization begins.

Completed model results record parameter count, convergence state, gradient
norm, regularization contribution, supported count mass, and, when supplied,
held-out likelihood and predictive error. Comparison utilities refuse to rank
models that lack held-out metrics; ranking is based on held-out evidence,
predictive error, and then parsimony, never on in-sample objective alone.
